\section{Accuracy-independent grids under quadratic growth}\label{sec:growth}

Theorem~\ref{thm:certificate} holds for any choice of grids. This section
chooses them so that, under weighted growth, every coordinate grid has
$O(\sqrt{\bar\kappa}\log(n+2))$ nodes at every refinement stage, independently
of the accuracy reached. The growth constant is used only in the analysis.

\subsection{Graded grids and one trial}

\begin{definition}[Graded grid]\label{def:graded}
Let $X_i'=[\ell_i',u_i']$ (or its integer points if $i\in I_Z$), a center
$c_i\in X_i'$, a mesh $h>0$ and a grading $0<\theta\le1/4$ be given. Put
\begin{equation}\label{eq:step}
\sigma(t)=h+\theta t\quad(i\in I_C),\qquad
\sigma(t)=\max\{1,\lfloor h+\theta t\rfloor\}\quad(i\in I_Z).
\end{equation}
Starting from $t_0=0$, set $t_{k+1}=\min\{t_k+\sigma(t_k),u_i'-c_i\}$ until
$t_{k+1}=u_i'-c_i$, and take the nodes $c_i+t_k$; proceed in the same way
towards $\ell_i'$. These nodes form the \emph{graded grid} of $X_i'$.
\end{definition}

Spacing grows linearly with the distance from the center, so a grid of radius
$R$ needs only $O(1+\theta^{-1}\log(1+\theta R/h))$ nodes.

\begin{lemma}[Graded grids]\label{lem:graded}
(a) $w_i(v)\le h+\theta\abs{v-c_i}$ for every node $v$.
(b) If an endpoint is at distance $R>0$ from $c_i$, the number of nodes
strictly on that side is at most $\lceil(4/\theta)\ln(1+\theta R/\hat h)\rceil$,
where $\hat h=h$ for $i\in I_C$ and $\hat h=\max\{h,1\}$ for $i\in I_Z$.
(c) If $h\ge u_i'-\ell_i'$, the grid has at most three nodes.
\end{lemma}

The proof is in Appendix~\ref{app:growth}. The meshes are chosen per
coordinate so that every coordinate sees the same correction scale. For
$i\in P$ let $\varpi_i\in\Z$ satisfy $4^{\varpi_i-1}<L_i\le4^{\varpi_i}$ and
put $r_i=2^{-\varpi_i}$, so that $\frac14<L_ir_i^2\le1$. Let $E$ be the least
integer with $2^Er_i\ge s_i$ for all $i\in P$, and put $\eta_j=2^{E-j}$ and
$h_{ij}=\eta_jr_i$; every $h_{ij}$ is a power of two.

\begin{algorithm}[TRIAL: one trial]\label{alg:trial}
The parameters are a grading $\theta\in(0,1/4]$, a cap $K$, an accuracy
$\varepsilon>0$ and a stage limit $j_{\max}$. Start with $X^{(0)}=X$, center
$c^{(0)}=\ell$, and the incumbent $\hat x=\ell$ with value $U=F(\ell)$, or a
better feasible point and its value if one is known. For
$j=0,1,\dots,j_{\max}$:
\begin{enumerate}[label=(\roman*)]
\item for $i\in P$ let $G_i^{(j)}$ be the graded grid of $X_i^{(j)}$ with
center $c_i^{(j)}$, mesh $h_{ij}$ and grading $\theta$; for $i\notin P$ let
$G_i^{(j)}=\{\ell_i,u_i\}$; if some grid would exceed $K$ nodes,
\emph{abort} the trial before forming any table;
\item compute $\beta_j$, a minimizer $y^{(j)}$ of $Q$ over $G^{(j)}$ and all
min-marginals by Lemma~\ref{lem:dp}; if $F(y^{(j)})<U$, set
$\hat x=y^{(j)}$ and $U=F(y^{(j)})$;
\item if $U-\beta_j\le\varepsilon$, stop with \emph{success}; otherwise
filter with threshold $U$ (Definition~\ref{def:filter}) to obtain
$X^{(j+1)}$, and set $c^{(j+1)}=y^{(j)}$.
\end{enumerate}
The trial \emph{fails} if stage $j_{\max}$ ends without success.
\end{algorithm}

By Proposition~\ref{prop:filter} the centers stay feasible and inside the
current subbox, and by Corollary~\ref{cor:filtercert} every stage yields a
valid path certificate. The analysis allows any first center in $X$
(Lemma~\ref{lem:inv}(ii) holds at $j=0$ for every $c^{(0)}\in X$). Later
centers need to be close to $x^*$: by Lemma~\ref{lem:inv}(iii), the corrected
minimizer $y^{(j)}$ is as close to $x^*$ as Lemma~\ref{lem:inv}(ii) requires
of the center of stage $j+1$. Because the centers are $\ell$ and grid points,
denominators do not compound from stage to stage
(Appendix~\ref{app:growth}).

\subsection{Contraction and localization}

\begin{lemma}[Stage invariants]\label{lem:inv}
Assume weighted growth with constant $\gamma$ at a minimizer $x^*$, let
$k\ge\max\{1,1/\gamma\}$ satisfy $8k\theta^2\le1$, and write
$a_j=n_P\eta_j^2$. At every stage $j$ reached by TRIAL, with $U$ the
incumbent value after step~(ii):
\begin{enumerate}[label=(\roman*)]
\item $x^*\in X^{(j)}$ and $\beta_j\le\OPT$;
\item $\norm{c^{(j)}-x^*}_{\mathbf L}^2\le4ka_j$;
\item $\norm{y^{(j)}-x^*}_{\mathbf L}^2\le ka_j$;
\item $U-\beta_j\le F(y^{(j)})-\beta_j=D(y^{(j)})\le\frac9{16}a_j$;
\item every grid point $z$ with $Q(z)\le U$ satisfies
$\norm{z-x^*}_{\mathbf L}^2\le\frac{17}{15}ka_j$.
\end{enumerate}
\end{lemma}

\begin{proof}
(i) By induction on $j$, starting from $x^*\in X^{(0)}=X$. If
$x^*\in X^{(j)}$, then $\beta_j\le\min_{X^{(j)}}F=\OPT$ by
Proposition~\ref{prop:cellwise}(b), and filtering with threshold
$U\ge\OPT=F(x^*)$ keeps $x^*$ by Proposition~\ref{prop:filter}.

Next, every grid point $z$ satisfies, with $c=c^{(j)}$,
\begin{equation}\label{eq:energy}
D(z)\le\frac{a_j}4+\frac{\theta^2}2\Bigl(\norm{z-x^*}_{\mathbf L}^2
+\norm{c-x^*}_{\mathbf L}^2\Bigr).
\end{equation}
Indeed, for $i\in P$ Lemma~\ref{lem:graded}(a) and $L_ih_{ij}^2\le\eta_j^2$
give $d_i(z_i)\le\frac{L_i}8(h_{ij}+\theta\abs{z_i-c_i})^2\le\frac{\eta_j^2}4
+\frac{\theta^2}4L_i(z_i-c_i)^2$, and
$(z_i-c_i)^2\le2(z_i-x_i^*)^2+2(c_i-x_i^*)^2$; for $i\notin P$, $d_i=0$.

(ii) For $j=0$ and any $c^{(0)}\in X$,
$\norm{c^{(0)}-x^*}_{\mathbf L}^2\le\sum_{i\in P}L_is_i^2\le
\sum_{i\in P}L_ir_i^2\eta_0^2\le a_0$. For $j\ge1$ it follows from (iii) at
stage $j-1$, because $c^{(j)}=y^{(j-1)}$ and $a_{j-1}=4a_j$.

(iii) Write $Y=\norm{y^{(j)}-x^*}_{\mathbf L}^2$ and $a=a_j$. Since
$\theta^2/2\le1/(16k)\le\gamma/16$, growth, (i), $Q(y^{(j)})=\beta_j$,
\eqref{eq:energy} and (ii) give
$\gamma Y\le F(y^{(j)})-\OPT\le F(y^{(j)})-\beta_j=D(y^{(j)})
\le\frac a4+\frac\gamma{16}(Y+4ka)$, so
$Y\le\frac4{15}(\gamma^{-1}+k)a\le\frac8{15}ka$.

(iv) By \eqref{eq:energy}, (ii), (iii) and $\theta^2k\le\frac18$,
$D(y^{(j)})\le\frac a4+\frac{\theta^2}2\cdot5ka\le\frac9{16}a$. Moreover
$U\le F(y^{(j)})$ after step~(ii), and $F(y^{(j)})-\beta_j=D(y^{(j)})$
because $Q(y^{(j)})=\beta_j$.

(v) Let $Z=\norm{z-x^*}_{\mathbf L}^2$. Using $Q(z)\le U\le F(y^{(j)})$,
$\beta_j\le\OPT$, (iv) and \eqref{eq:energy},
$\gamma Z\le F(z)-\OPT=Q(z)+D(z)-\OPT\le D(y^{(j)})+D(z)\le\frac{13}{16}a
+\frac\gamma{16}(Z+4ka)$, so
$Z\le(\frac{13}{15}\gamma^{-1}+\frac4{15}k)a\le\frac{17}{15}ka$.
\end{proof}

The proof shows where filtering is essential. Part (v) bounds every grid point
whose corrected value can keep an interval alive, and the retained subbox
therefore shrinks at the same rate as the mesh.

\begin{lemma}[Localization and grid size]\label{lem:states}
Under the hypotheses of Lemma~\ref{lem:inv}, for every stage $j$ followed by
filtering and every $i\in P$,
\[
X_i^{(j+1)}\subseteq\bigl[y_i^{(j)}-R_{ij},\,y_i^{(j)}+R_{ij}\bigr],\qquad
R_{ij}=8\sqrt{kn_P}\,h_{ij}+[i\in I_Z],
\]
where $[i\in I_Z]$ is $1$ if $i\in I_Z$ and $0$ otherwise. Consequently
every grid of the trial has at most
\begin{equation}\label{eq:statecap}
K(\theta,n_P)=10\,\theta^{-1}\lceil\log_2(n_P+2)\rceil
\end{equation}
nodes per coordinate, independently of the stage and of $\varepsilon$.
\end{lemma}

\begin{proof}
Write $y=y^{(j)}$, $c=c^{(j)}$, $h=h_{ij}$ and $\rho=\sqrt{kn_P}$. Let
$[a,a']$ be a retained grid interval of coordinate $i$, let
$v\in\{a,a'\}$ satisfy $m_i(v)\le U$ (Definition~\ref{def:filter}), and let
the grid point $z$ attain $m_i(v)$, so $z_i=v$ and $Q(z)\le U$. Because
$\sqrt{L_i}>1/(2r_i)$, a bound $\sqrt{L_i}\abs\delta\le\lambda\rho\eta_j$
implies $\abs\delta<2\lambda\rho h$; so Lemma~\ref{lem:inv}(ii), (iii), (v)
give $\abs{v-y_i}<2(1+\sqrt{17/15})\rho h$ and
$\abs{v-c_i}<2(2+\sqrt{17/15})\rho h$. An integer interval of length one
lies within $\abs{v-y_i}+1$ of $y_i$. Otherwise the interval has length at
most $w_i(v)\le h+\theta\abs{v-c_i}$ (Lemma~\ref{lem:graded}(a)), and
$\theta\sqrt k\le1/\sqrt8$, $h\le\rho h$ and $\sqrt{n_P}\le\rho$ give that
each of its points is within
$\bigl(3+2\sqrt{17/15}+(2+\sqrt{17/15})/\sqrt2\bigr)\rho h<7.3\rho h$ of
$y_i$. So every retained interval, and hence their hull $X_i^{(j+1)}$, lies
in $[y_i-R_{ij},y_i+R_{ij}]$.

At stage $j+1$ the grid of coordinate $i$ is centered at
$y_i\in X_i^{(j+1)}$ with mesh $h/2$, and each side has length at most
$R_{ij}$. With $\hat h$ as in Lemma~\ref{lem:graded}(b) for mesh $h/2$,
$\theta R_{ij}/\hat h\le16\theta\rho+\theta\le4\sqrt{2n_P}+\frac14$, so by
Lemma~\ref{lem:graded}(b) the grid has at most
$1+2\lceil\frac4\theta\ln(\frac54+4\sqrt{2n_P})\rceil$ nodes, which is at
most $K(\theta,n_P)$ by a calculus estimate in Appendix~\ref{app:growth}.
Stage $0$ has at most three nodes per
coordinate by Lemma~\ref{lem:graded}(c), since $h_{i0}\ge s_i$, and
coordinates outside $P$ have two.
\end{proof}

Section~\ref{sec:related} relates Lemmas~\ref{lem:inv} and~\ref{lem:states}
to the cluster problem of branch-and-bound.

\subsection{Unknown growth: capped trials}
The algorithm cannot use $8\bar\kappa\theta^2\le1$ directly because $\gamma$
is unknown. It tries geometrically decreasing gradings and relies on the caps
to bound the work of unsuitable trials.

\begin{algorithm}[CT: capped trials]\label{alg:ct}
The input is an accuracy $\varepsilon>0$. If $P=\emptyset$, minimize $F$
over the grid $\prod_i\{\ell_i,u_i\}$ by Lemma~\ref{lem:dp} and return a
minimizer $\hat x$, $\beta=F(\hat x)$ and this grid. Otherwise let
$j_{\max}$ be the least integer $j_{\max}\ge0$ with
$\frac9{16}n_P\eta_0^2\,4^{-j_{\max}}\le\varepsilon$, and for
$\mu=2,3,\dots$ run TRIAL with grading $\theta=2^{-\mu}$, cap
$K_\mu=K(2^{-\mu},n_P)$, accuracy $\varepsilon$ and stage limit $j_{\max}$,
keeping the incumbent from trial to trial, until a trial succeeds. Return the
incumbent $\hat x$, the bound $\beta=\beta_j$ of the successful stage $j$ and
the grids of the successful trial.
\end{algorithm}

\begin{theorem}[Certified approximation]\label{thm:approx}
Let~\eqref{eq:model} have explicit rational quadratic factors given as
monomial lists, a tree decomposition of maximum bag size $p$, curvature bounds
$L_i=\max\{0,\partial_{ii}F\}$, and let $\varepsilon=2^{-q}$.
\begin{enumerate}[label=(\alph*)]
\item On every instance, CT$(\varepsilon)$ halts and returns a feasible
rational $\hat x$ and a valid path certificate with
$\beta\le\OPT\le F(\hat x)\le\beta+\varepsilon$.
\item If the instance has weighted growth with condition number $\bar\kappa$,
CT$(\varepsilon)$ succeeds at a trial with $2^\mu\le6\sqrt{\bar\kappa}$ and
uses at most
\[
f(p,\bar\kappa)\,(I+q+1)^5,\qquad
f(p,t)=c_0\,(c_1p\sqrt t)^p\,t\,(1+\log_2t)^2,
\]
bit operations, for absolute constants $c_0,c_1$. The certificate has size,
and is checked in time, within the same bound. Neither $\gamma$ nor
$\bar\kappa$ is an input. Since $f$ is increasing in $t\ge1$, the bound also
holds with any larger number in place of $\bar\kappa$, such as the Euclidean
condition number $\kappa$ under point growth.
\end{enumerate}
\end{theorem}

\begin{proof}
(a) If $P=\emptyset$, all corrections vanish and
Proposition~\ref{prop:cellwise}(b) gives $F(\hat x)=\beta\le\OPT$. Otherwise
validity follows from Corollary~\ref{cor:filtercert}, because every threshold
is the value of a feasible point. For termination, take $\mu\ge2$ with
$\theta=2^{-\mu}\le\min_{i\in P}h_{i,j_{\max}}/s_i$, so that
$\theta s_i\le h_{ij}$ for all $j\le j_{\max}$. In this trial every unclipped
step at a stage $j\le j_{\max}$ is at least $h_{ij}$: for $i\in I_C$
by~\eqref{eq:step}, and for $i\in I_Z$ with
$h_{ij}\ge1$ because $h_{ij}$ is a power of two, so
$\lfloor h_{ij}\rfloor=h_{ij}$. A side of length at most $s_i\le h_{ij}/\theta$
therefore takes at most $1/\theta$ steps, and the grid has at most
$1+2/\theta<K_\mu$ nodes. For $i\in I_Z$ with $h_{ij}<1$ the grid has at most
$s_i+1<1/\theta+1$ nodes, because $s_i\le h_{ij}/\theta<1/\theta$. So the
trial is not aborted. At stage $j_{\max}$, Lemma~\ref{lem:graded}(a) gives
$w_i(v)\le h_{i,j_{\max}}+\theta s_i\le2h_{i,j_{\max}}$ for every node $v$, so
$U-\beta_{j_{\max}}\le D(y^{(j_{\max})})\le\sum_{i\in P}L_ih_{i,j_{\max}}^2/2
\le a_{j_{\max}}/2\le\frac89\varepsilon$, and the trial succeeds.

(b) Let $\mu^*=\max\{2,\lceil\log_4(8\bar\kappa)\rceil\}$, so that
$\theta=2^{-\mu^*}$ satisfies $8\bar\kappa\theta^2\le1$, and
$2^{\mu^*}\le6\sqrt{\bar\kappa}$. Apply Lemmas~\ref{lem:inv}
and~\ref{lem:states} with the growth constant $\gamma$ that defines
$\bar\kappa$ and $k=\bar\kappa$: by Lemma~\ref{lem:states} trial $\mu^*$ is
never aborted, and by Lemma~\ref{lem:inv}(iv) it succeeds at stage
$j_{\max}$ at the latest, since $\frac9{16}a_{j_{\max}}\le\varepsilon$. Every
trial $\mu\le\mu^*$ runs at most $j_{\max}+1$ stages with at most $K_\mu$
nodes per coordinate, because an oversized grid is detected while it is
generated. By Lemma~\ref{lem:dp} this costs
$O\bigl((j_{\max}+1)p(\abs{\mathcal A}+n+N)K_\mu^p\bigr)$ table operations,
and $\sum_{\mu\le\mu^*}K_\mu^p\le2K_{\mu^*}^p$. Since
$\sup_{t\ge0}t^pe^{-t}=(p/e)^p$,
\begin{equation}\label{eq:logabsorb}
\lceil\log_2(n_P+2)\rceil^p\le2p^p(n+2),
\end{equation}
so the table work is at most
$c\,(j_{\max}+1)\,p\,I^2(60\,p\sqrt{\bar\kappa})^p$ with
$j_{\max}=O(I+q+1)$ and an absolute constant $c$.
Let $\Gamma_X$ and $\Gamma_F$ be the products of the denominators of the box
endpoints and of the coefficients. Appendix~\ref{app:growth} shows that all
grid nodes of a trial have denominators dividing $\Gamma_X2^\alpha$, where
$\alpha=j_{\max}+O(I)+\mu K_\mu$, so one bound holds at all stages of a
trial: denominators do not compound from stage to stage. It also shows that
all table entries are integers of $O((1+\mu2^\mu)(I+q+1))$ bits over the
common denominator $8\Gamma_F\Gamma_X^24^\alpha$. Multiplying the counts
proves the bound.
\end{proof}

\begin{remark}[In what sense the bound is fixed-parameter]\label{rem:fpt}
As a parameter, $\bar\kappa$ is the smallest weighted condition number of
the instance, which exists by Section~\ref{sec:setting}. It is a real number,
not part of the input, and CT never computes it. Since $f$ is
increasing in $t\ge1$, the bound of Theorem~\ref{thm:approx}(b) is at most
$f(p,\lceil\bar\kappa\rceil)(I+q+1)^5$, a fixed-parameter bound in the
parameters $p$ and $\lceil\bar\kappa\rceil$, which the algorithm does not
need to know. The exponent $5$ reflects schoolbook arithmetic and depends on
neither parameter, and the accuracy enters only through
$q=\log_2(1/\varepsilon)$. The factor $\lceil\log_2(n_P+2)\rceil^p$ is
absorbed by~\eqref{eq:logabsorb}, and no bound on the number of bags
containing a coordinate, on Hessian norms or on the number of local minima is
used. If only a treewidth bound is known, the decomposition of
\cite{Korhonen2021} gives the same result with $p=2\,\mathrm{tw}+2$.
Theorem~\ref{thm:intro-lower} states the lower bounds, proved in
Section~\ref{sec:limits}; since $\bar\kappa\le\kappa$, the bounds stated
for $\kappa$ also apply to $\bar\kappa$.
\end{remark}

\paragraph{A common mesh.}
The implementation of Section~\ref{sec:computation} and the exact-output
results of Section~\ref{sec:localized} and
Appendices~\ref{app:localized} and~\ref{app:boundary} use one mesh for all
coordinates and point growth with the Euclidean condition number $\kappa$.

\begin{lemma}[CT with a common mesh]\label{lem:commonmesh}
Modify TRIAL by using the common mesh $h_j=s2^{-j}$ in place of $h_{ij}$ for
every $i\in P$; the corrections $L_iw_i(v)^2/8$ and the grids
$\{\ell_i,u_i\}$ for $i\notin P$ are unchanged. Assume point growth with
constant $g$ at $x^*$, put $\kappa=\max\{1,L/g\}$, and let
$8\kappa\theta^2\le1$. At every stage $j$ reached by such a trial, with
center $c=c^{(j)}$, corrected minimizer $y_j=y^{(j)}$ and $U$ the incumbent
value after step~(ii), we have $x^*\in X^{(j)}$, $\beta_j\le\OPT$, and:
\begin{enumerate}[label=(G\arabic*)]
\item $D(z)\le\frac L4nh_j^2+\frac{L\theta^2}2\bigl(\norm{z-x^*}^2
+\norm{c-x^*}^2\bigr)$ for every grid point $z$;
\item $\norm{c-x^*}^2\le4\kappa nh_j^2$, $\norm{y_j-x^*}^2\le\kappa nh_j^2$,
and $\norm{z-x^*}^2\le\frac{17}{15}\kappa nh_j^2$ for every grid point $z$
with $Q(z)\le U$;
\item $U-\OPT\le F(y_j)-\OPT\le F(y_j)-\beta_j=D(y_j)\le\frac9{16}Lnh_j^2$;
\item if stage $j$ is followed by filtering, then for $i\in P$ the filtered
interval $X_i^{(j+1)}$ lies within $4.2\sqrt{n\kappa}\,h_j$ of $(y_j)_i$ if
$i\in I_C$, and within $1+4.2\sqrt{n\kappa}\,h_j$ if $i\in I_Z$;
\item every grid of the trial has at most
$8\theta^{-1}\lceil\log_2(n+2)\rceil$ nodes per coordinate.
\end{enumerate}
These statements also hold if the trial starts at any center
$c^{(0)}\in X$ instead of $\ell$.
\emph{CT with a common mesh} is CT with these trials, caps
$8\cdot2^\mu\lceil\log_2(n+2)\rceil$, and the least stage limit
$j_{\max}\ge0$ with $\frac9{16}Lns^24^{-j_{\max}}\le\varepsilon$. On every
instance it halts with a valid path certificate and
$F(\hat x)\le\beta+\varepsilon$; under point growth it succeeds at the latest
in the trial
$\mu=\max\{2,\lceil\log_4(8\kappa)\rceil\}$.
\end{lemma}

The proof in Appendix~\ref{app:growth} follows those of
Lemmas~\ref{lem:inv} and~\ref{lem:states} with $a_j=Lnh_j^2$ and
$\norm\cdot_{\mathbf L}^2$ replaced by $L\norm\cdot^2$. Because (G2) bounds
Euclidean distances in units of the common mesh, the coordinate bounds
follow directly, without the conversion through $\sqrt{L_i}>1/(2r_i)$ that
costs Lemma~\ref{lem:states} up to a factor~$2$; this gives the smaller
radius and cap, at the price of $\kappa\ge\bar\kappa$. Items (G1)--(G4) and
their proofs also hold for uniform grids, $\theta=0$, with
Lemma~\ref{lem:graded}(a) read as $w_i(v)\le h_j$; only (G5) needs
$\theta>0$.

\subsection{Sharpness of the localization radius}
The slack $U-\beta$ can be as large as a sum of $n$ correction terms,
while a single coordinate interval is removed only when its own excess
exceeds the whole slack. The following family shows that the radius
$\Theta(\sqrt{n\kappa}\,h_j)$ of Lemmas~\ref{lem:states}
and~\ref{lem:commonmesh} is therefore attained, so that filtering alone
leaves at least of order $\sqrt{n\kappa}$ nodes per coordinate on uniform
grids; the graded grids reduce this to
$O(\sqrt{\bar\kappa}\log(n+2))$ (Lemma~\ref{lem:states}), which makes
Theorem~\ref{thm:approx} a fixed-parameter bound (Remark~\ref{rem:fpt}).

\begin{proposition}[Sharpness of localization]\label{prop:sharp}
Let $m\ge1$, $n=2m$, $\Lambda>0$ and $0<g\le\Lambda/2$, and on
$X=[-1,1]^n$, with all coordinates continuous and named
$(u_1,v_1,\dots,u_m,v_m)$, let
\[
F(x)=\sum_{k=1}^m\Bigl[\frac\Lambda2(u_k^2+v_k^2)+(\Lambda-2g)u_kv_k\Bigr].
\]
Then $x^*=0$, $\OPT=0$, the smallest valid curvature bounds are
$L_i=\Lambda$, the largest growth constant is $g$, and $\kappa=\Lambda/g$
takes every value in $[2,\infty)$. Let $G$ be a grid for a subbox that
contains $[-R_0,R_0]^n$, and let every $G_i$ contain $0$ with
$w_i(0)\ge h$. Then for every incumbent value $U\ge0$, every coordinate $i$
and every node $a\in G_i$ with
\[
|a|\le\min\Bigl\{R_0,\ h\sqrt{(n-2)\kappa/16}\Bigr\}
\]
we have $m_i(a)\le U$. Hence every grid interval with such an endpoint is
retained.
\end{proposition}

\begin{corollary}[Filtered uniform grids need order $\sqrt{n\kappa}$ nodes]
\label{cor:uniformgrid}
Take the instance of Proposition~\ref{prop:sharp} with $n\ge4$, and run the
stages of a trial with a common mesh (Lemma~\ref{lem:commonmesh}) without a
cap, with $h_j=s2^{-j}=2^{1-j}$, uniform grids (the construction of
Definition~\ref{def:graded} with $\theta=0$) and initial center $0$. Let
$k_0=\lfloor\sqrt{(n-2)\kappa/16}\rfloor$. At every stage the corrected
minimizer is $0$; at every stage $j\ge1$ with $(k_0+1)h_j\le1$ that ends by
filtering, the filtered box $X^{(j+1)}$ contains
$[-(k_0+1)h_j,(k_0+1)h_j]^n$, and the grid of stage $j+1$, if that stage is
run, has at least $4k_0+5\ge\sqrt{(n-2)\kappa}+1$ nodes in every coordinate.
With graded grids of grading $\theta\in(0,1/4]$ instead: if a stage $j$ ends
by filtering, its box $X^{(j)}$ contains $[-R,R]^n$, where
$R=h_j\sqrt{(n-2)\kappa/16}$, and its grids contain $0$ with
$w_i(0)\ge h_j$, then the grid of stage $j+1$, if that stage is run, has at
least $1+\ln(1+2\theta R/h_j)/\ln(1+\theta)$ nodes in every coordinate.
\end{corollary}

Proofs are in Appendix~\ref{app:growth}. Uniform grids also give
accuracy-independent counts without any grading condition: under point
growth, UC (Algorithm~\ref{alg:uc}), which filters unions of uniform cells,
has at most $12(2\sqrt{n\kappa}+1)$ nodes per coordinate at every stage
(Theorem~\ref{thm:cells} with $r=1$ and $\kappa_S=\kappa$). The resulting
running-time bound is polynomial in $I+q$ and $\kappa$ for fixed $p$, but
because of the factor $\sqrt n$ per coordinate it is not a fixed-parameter
bound.


\subsection{Examples}
The first example shows that Theorem~\ref{thm:approx} covers nonconvex
problems with many local minima and treewidth larger than one.

\begin{example}[Coupled indefinite blocks]\label{ex:family}
For a block $b$ with variables $(u_b,v_b,\omega_b)\in[0,1]^3$ let
$\phi(u,v,\omega)=u^2+v^2-4uv+\frac14(u+v)+(\omega-\frac u2)^2$, and for a graph
$\mathcal G$ on the blocks with maximum degree $\Delta_{\mathcal G}\ge1$ let
\[
F=\sum_b\phi(u_b,v_b,\omega_b)+\frac1{16\Delta_{\mathcal G}}
\sum_{bb'\in\mathcal G}(u_b-u_{b'})^2 .
\]
Then (a) $x^*=(1,1,\frac12)$ in every block is the unique minimizer and
$F(x)-\OPT\ge\frac12\norm{x-x^*}^2$; (b) $L\le\frac{21}8$, so
$\kappa\le\frac{21}4$; (c) each of the $2^m$ points with
$(u_b,v_b,\omega_b)\in\{(0,0,0),(1,1,\frac12)\}$ in all $m$ blocks is a strict
local minimizer; (d) a tree decomposition of $\mathcal G$ of width $w$ yields
one of $F$ with $p=\max\{w+1,3\}$. Each block Hessian is indefinite and the
minimizer is not a vertex. For $\mathcal G$ an $m_1\times m_2$ grid graph
with $m=m_1m_2\ge2$, Theorem~\ref{thm:approx} applies with $p=\max\{m_1+1,3\}$ and
$\bar\kappa\le\kappa\le\frac{21}4$, although the instance has $2^m$ strict
local minima. The proof is in Appendix~\ref{app:growth}.
\end{example}

The minimizer of this family is easy to guess block by block; by (a),
$F-\OPT$ is at least $\frac98$ times the number of blocks at $(0,0,0)$. The
second example separates the grids from exact parametric dynamic
programming.

\begin{example}[Expanding-box chain]\label{ex:chain}
The family of Proposition~\ref{lim:prop:messages}, a chain of length $m$
whose $t$th state has a box of width $2^t-1$, has $p=3$, a unique minimizer
and $\kappa\le80$ for every $m$. By Theorem~\ref{thm:approx}, a certified
$2^{-q}$-approximation takes time polynomial in $m$ and $q$, whereas the
exact Bellman message for the last state has at least $2^{m-1}$ quadratic
pieces. Rescaling the $t$th state by $2^{-t}$ moves every coordinate into
$[0,1]$ and does not change $\bar\kappa$, which stays at most $80$, but the
rescaled instance has $L=4^{m+1}$ and a largest growth constant between
$\frac18$ and $20$, so every valid growth constant gives $\kappa\ge4^m/5$,
and the largest gives $\kappa\le32\cdot4^m$ (Appendix~\ref{app:growth}).
Section~\ref{sec:computation} reports the certified computation up to
$m=64$.
\end{example}
